\section{R1: Search control with coupled goals}\label{sec:r1}

\question{R1.Q1}{Can shared-metavariable goals be organized without wholesale copying or sequentialization?}
\status{Yes, through conditional factorization and tabling; no, if the demand is to erase the information distinguishing incompatible assignments.}

Represent a node as a constraint hypergraph. Its vertices include unfinished
term holes, universe variables, postponed constraints, and observations.
A hyperedge records a joint dependency. Dependencies must be closed through
assigned expressions, delayed assignments, local declaration types and values,
and instance arguments. For example, two goals whose printed targets differ
may depend on the same universe variable or on a local whose type contains a
shared hole. Direct comparison of target metavariable sets misses this.

A component can be solved independently only after its external interface is
fixed, or while its solutions remain parameterized by that interface. Store a
solution as an abstracted term and a constraint on its interface, not as an
unqualified assignment to a global metavariable identifier. On reuse,
instantiate fresh variables, check the interface constraints, and replay the
term through type checking. For small separators this resembles variable
elimination: each side exports a relation on a few shared variables rather
than every internal branch.

\begin{proposition}[A sufficient factorization criterion]\label{prop:factor}
Suppose a state's completion constraints are $D_1\land D_2$, where the
unassigned variables of $D_1$ and $D_2$ are disjoint after dependency closure.
Assume a fixed common telescope and that every allowed action on either side
preserves this separation. Then compatible solutions of $D_1$ and $D_2$
combine by disjoint substitution, and every joint solution restricts to a
solution of each side.
\end{proposition}
\begin{proof}
The substitutions have disjoint domains and no right-hand-side dependence on
the other side's unassigned variables. Their union is therefore a well-defined,
capture-avoiding substitution. Substitution into $D_i$ is unaffected by the
other component, so both constraints hold. Restriction gives the converse.
The action-preservation condition prevents later creation of a hidden shared
constraint from invalidating the factorization.
\end{proof}

This is a sufficient condition, not a cheap complete independence test.
Start with syntactic dependency closure and conservatively retain uncertain
edges. Recompute only after relevant assignments or new observations. Cache
successful, abstracted solutions first; caching failure requires a certificate
of exhaustion at a particular grammar and resource bound.

Persistent contexts already share most unchanged structure. Aesop already
computes transitive clusters of coupled goals, including dependencies through
local contexts. Its copying concerns goal nodes, not wholesale copies of their
search subtrees. Thus cheap independence detection is not a missing idea in
Aesop; the additional proposal here is interface-parameterized solution
tabling and carefully limited factorization \cite{aesop}. Incompatible
assignments still require distinct logical branches.

There is also an immediate implementation caveat: \code{Meta.SavedState} is
not the whole current search state. The residual-blocker array and several
other fields are \code{IO.Ref}s outside the snapshot. A queued node needs its
own residual-cache identity and resumable expansion cursor. Global work
charges should remain outside snapshots, but branch-dependent semantic caches
must not be shared without valid keys \cite{searchcode,transactioncode}.

\question{R1.Q2}{What admissible heuristic remains valid when unification closes several goals?}
\status{Use zero as the baseline, maximum rather than sum for coupled goals, and additive bounds only across action-separable components.}

Suppose two unfinished obligations are a data hole $x:\Nat$ and an equality
hole $p:x=7$. An equality-solving step may instantiate $x$ while solving $p$.
If that step costs one, charging one for each obligation gives an estimate of
two for a remaining cost of one. Even an isolated goal may have a zero-cost
exact solution under the proposal's weights. The number of visible goals is
therefore not automatically a lower bound.

\begin{proposition}[Combining lower bounds]\label{prop:heuristic}
Let $c^*(s)$ be the minimum additional cost of completing state $s$, with
nonnegative action costs. If $h_i(s)\le c^*(s)$ for every $i$, then
$\max_i h_i(s)\le c^*(s)$. A sum of component bounds is admissible if every
completion's cost can be partitioned into component costs, and the bound for
each component is no greater than its corresponding share.
\end{proposition}
\begin{proof}
The maximum cannot exceed a common upper bound. Under the partition
hypothesis, sum the inequalities for the component shares. Without that
hypothesis, one action can pay for several components and the sum double-counts.
\end{proof}

A better nontrivial bound comes from a finite abstraction of the remaining
constraint system. Give every concrete action an abstract counterpart of no
greater cost; allow the abstraction to forget constraints and to close several
abstract goals together. Abstract shortest-path distance is then admissible:
every concrete completion projects to an abstract completion of no larger
cost. An abstraction that merely treats all goals independently does not have
this property without an additional proof.

Use a learned estimate of difficulty for ordering, but do not call it an
admissible heuristic. Keep an immutable logical cost for optimality claims.
Also, zero-cost rules need either an acyclic normalization phase, duplicate
suppression with a termination argument, or a structural bound: infinitely
many zero-cost expansions can otherwise prevent progress in cost order.

\question{R1.Q3}{Can just-in-time learning preserve meaningful ranking?}
\status{It preserves logical soundness under the gate, but neither budget invariance nor an immutable ranking objective follows automatically.}

Probe is an important precedent for using partial success to guide enumeration,
not a theorem that learned search order is independent of the examples,
execution trace, or stopping time \cite{probe}. Maintain two quantities:
\[
 c_0(t)=\text{fixed description cost},\qquad
 q_e(s)=\text{exploration priority during epoch }e.
\]
Use $q_e$ to allocate search effort and $c_0$ to compare completed candidates.
A learning update must occur after a deterministic logical event, such as a
fixed number of completed expansions, not after a wall-clock interval or a
race between parallel workers. Record the update and its inputs.

For best-first search with changing priorities, either rebuild affected queue
keys or freeze priorities for a bounded epoch. An A* proof about old keys is
not valid after arbitrary weight changes. A fixed-cost companion frontier can
supply a lower-bound certificate even while a heuristic lane explores other
nodes. A reserved fair share for that companion prevents learned priorities
from starving a valid bounded alternative.

Users should see ``best found under this search receipt,'' unless the fixed
frontier proves a stronger statement. A reference solution moving in the list
because the algorithm found better evidence is defensible; pretending that an
adaptive search returns a machine-independent optimum within ten seconds is
not. Learning may use the proportion of observations satisfied, but the
observations should be weighted carefully: forty correlated tests should not
necessarily overwhelm one structurally different test.

\question{R1.Q4}{Can a hard wall-clock budget give machine-independent results?}
\status{Not in general. Deterministic prefixes, replay, and completed logical epochs are achievable alternatives.}

\begin{proposition}[A wall-clock obstruction]\label{prop:clock}
Consider a deterministic algorithm whose answers change after some positive
amount of computation. If execution speed may vary sufficiently, a common hard
wall-clock cutoff cannot guarantee the same answer set on every execution
while also publishing all answers reached before that cutoff.
\end{proposition}
\begin{proof}
Choose a logical event at which the answer set changes. Choose two speeds so
that the event occurs before the deadline on one execution and after it on the
other. Publishing the answers reached gives different sets. Suppressing the
new answer avoids that pair of outputs only by abandoning the requirement to
publish all reached answers.
\end{proof}

Even a restricted speed variation breaks equality whenever a cutoff straddles
such an event. A known worst-case execution bound can support a fixed amount
of logical work, but then the guarantee is conditional on that bound. Startup
calibration is an estimate, not a proof; cache warmth and input-dependent
kernel costs can invalidate it.

The practical contract should be \emph{prefix reproducibility}: pin the
query, environment, grammar, ranking model, and tie-breaking; number logical
expansions; publish a resumable trace prefix. A deterministic-work run can
reproduce that prefix regardless of elapsed time. A timed run returns the
prefix it managed to finish. Completed epochs may be compared across machines,
but their number need not match. Parallel workers may explore speculatively,
while a deterministic merger commits results in logical order.

Luby-style restart schedules give useful guidance under their Las Vegas
runtime assumptions; arbitrary coupled, stateful lanes do not automatically
satisfy those assumptions \cite{luby}. Prefer resumable lanes and deterministic
work allocation before adding restarts. Replaying the same initial search
because a deadline interrupted it is wasted work unless the restart changes
something that justifies its cost.
